% 研究方法：共同設定集中說明，實驗依執行順序編排。
\chapter{研究方法}\label{chap:method}
\begingroup
\setlength{\emergencystretch}{2em}
\clubpenalty=10000
\widowpenalty=10000
\displaywidowpenalty=10000

\section{研究架構}\label{sec:method-framework}

本章說明兩階段集成式持續學習的操作流程。實驗一先比較單一分類器的訓練期間與採樣方式，再建立 Old／New 模型池，比較靜態集成、DCS 與 DES；實驗二在相同條件下加入特徵選取，以完整特徵結果為對照。整體流程如圖~\ref{fig:method-overview}，實際資料與年份配置見第~\ref{sec:method-data}~節，評估與統計程序見第~\ref{sec:method-validation}~節。

輸入資料表示為 \(\{(x_i,y_i,t_i)\}_{i=1}^{n}\)，其中 \(x_i\) 為特徵向量、\(y_i\in\{0,1\}\) 為分類標籤、\(t_i\) 為觀測時間，\(n\) 為樣本數。同一個體的識別碼另供分區檢核與配對。以 \(\mathcal{O}\) 表示歷史模型池、\(\mathcal{N}\) 表示近期新增模型，更新由 \(\mathcal{O}\) 擴充至 \(\mathcal{O}\cup\mathcal{N}\)，再以驗證資料決定使用成員與分類門檻。本次採固定時間分區，在同一離線實驗內重建各設定，模擬一次模型池更新。

\begin{figure}[htbp]
\centering
\begingroup
\small
\setstretch{1.12}
\setlength{\fboxsep}{8pt}
\begin{tabular}{@{}c@{}}
\fbox{\parbox{0.88\linewidth}{\centering 具有時間資訊的分類資料\\
定義觀測單位與正類；依時間劃分擬合、驗證與測試用途}}\\[3pt]
\(\Downarrow\)\\[3pt]
\fbox{\parbox{0.88\linewidth}{\centering\textbf{實驗一：Baseline 與集成分類器比較}\\[4pt]
Baseline：分類器、Old／New／Retrain、採樣策略之比較\\
\(\Downarrow\)\\
以驗證選定的最佳模型建立 Old 三模型與 New 三模型\\
比較靜態集成、DCS 與 DES}}\\[3pt]
\(\Downarrow\)\\[3pt]
\fbox{\parbox{0.88\linewidth}{\centering\textbf{實驗二：特徵選取集成分類器比較}\\[4pt]
沿用集成比較之資料分區、分類器種類與採樣方式\\
以 Old 擬合資料選取特徵，再採樣並重新建立六模型池\\
比較加入 FS 前後的固定方法與驗證選模程序}}\\[3pt]
\(\Downarrow\)\\[3pt]
\fbox{\parbox{0.88\linewidth}{\centering 驗證決策固定後，評估保留的較晚期間測試資料\\
主要報告 AUC、F1-score 與 Recall}}
\end{tabular}
\endgroup
\caption{本研究兩項實驗之整體架構}\label{fig:method-overview}
\par\smallskip{\footnotesize (資料來源：本研究)\par}
\end{figure}

\clearpage

\section{實驗資料與配置}\label{sec:method-data}

本節交代兩項實驗共用的資料、時間分區與處理規則。各方法的原理已於第二章說明，以下著重實際輸入與執行設定。

\subsection{資料來源與標籤定義}\label{sec:method-label}

本研究使用 Lombardo 等人\autocite{lombardo2022bankruptcy}所提供、於 2024 年 4 月 10 日更新的公開資料\autocite{sowide-bankruptcy-data}，包含 1999--2018 年的 78,682 筆公司年度紀錄、8,971 個匿名公司識別碼與 18 個財務特徵。公開說明的 8,262 家公司與檔內識別碼數不同，本文依實際檔案計數；版本核對見附錄~\ref{sec:appendix-data-fs}。

原始欄位 \nolinkurl{status_label} 記錄公司的最終狀態：609 個識別碼為 \nolinkurl{failed}，其所有年度共 5,220 筆紀錄均標為破產，占 6.63\%。主要分析改用回溯重建的事件代理標籤：破產公司僅最後觀測年度記為 1，作為推定下一年度事件的代理；其餘紀錄記為 0，共有 609 筆代理事件，占 0.77\%。年度 \(t\) 的標籤可用時間設為 \(t+1\)，並據此檢查各分區標籤在測試開始時是否成熟。此為年度可用性假設；公開檔缺乏逐筆破產申請日期與財報公開時間，無法驗證實際資訊到達日。

最終狀態標籤保留為對照，使用相同資料列、分區及模型池結構，只改變目標欄位；結果見第~\ref{sec:results-label}~節。

\subsection{輸入欄位與資料品質}

觀測單位為公司年度紀錄，模型輸入限於表~\ref{tab:method-feature-dictionary} 的 X1--X18。識別碼、年度、標籤及產業分類均不作預測特徵，亦未新增比率、股價或落後期特徵。同一公司可跨擬合與測試年份出現，因此評估對象是較晚年度的紀錄，並非完全未見的公司。

載入時檢查公司與年度組合及財務欄位有效值，未發現重複公司年度紀錄、遺失值或無限值，故沒有因此刪除紀錄。各設定仍保留相同的補值器，完整處理規則見第~\ref{sec:method-common}~節。

2024 年公開版的欄位順序及數值與 2022 年版不同，本研究依公開研究附件的欄位名稱及跨版本數值對照整理\autocite{sowide-financial-supplement}，保留實際版本的數值與順序，未依舊版倍率換算。

\Needspace{21\baselineskip}
\begingroup
\small\setstretch{1.12}\setlength{\tabcolsep}{4pt}
\begin{longtable}{@{}>{\raggedright\arraybackslash}p{0.15\linewidth}>{\raggedright\arraybackslash}p{0.77\linewidth}@{}}
\caption{本研究使用之財務特徵}\label{tab:method-feature-dictionary}\\
\toprule
欄位 & 財務項目 \\
\midrule\endfirsthead
\multicolumn{2}{c}{表~\thetable（續）}\\\toprule
欄位 & 財務項目 \\
\midrule\endhead
\midrule\multicolumn{2}{r}{續下頁}\\\endfoot
\bottomrule\endlastfoot
X1 & 流動資產 \\
X2 & 資產總額 \\
X3 & 銷貨成本 \\
X4 & 長期負債總額 \\
X5 & 折舊與攤銷 \\
X6 & 稅息前盈餘（EBIT） \\
X7 & 稅息折舊攤銷前盈餘（EBITDA） \\
X8 & 毛利 \\
X9 & 存貨 \\
X10 & 流動負債總額 \\
X11 & 淨利 \\
X12 & 保留盈餘 \\
X13 & 應收款項總額 \\
X14 & 營收總額 \\
X15 & 市值 \\
X16 & 負債總額 \\
X17 & 銷貨淨額 \\
X18 & 營業費用總額 \\
\end{longtable}
\noindent{\footnotesize (資料來源：本研究依作者公開研究附件之欄位名稱與跨版本資料對照整理\autocite{sowide-financial-supplement})\par 本表使用本研究實際版本的順序；X14 與 X17 在目前公開檔中的數值相同，名稱依發布附件保留。}
\endgroup

\subsection{時間切分與驗證用途}\label{sec:method-time-window}

固定年份分區列於表~\ref{tab:method-partitions}。歷史／近期的切點及測試區間參考原資料研究\autocite{lombardo2022bankruptcy}，但另將 2014 年完整保留供驗證使用，使 New 實際只擬合 2012--2013 年。由於可用訓練資料與原文不同，不直接比較兩者分數；2011 年切點亦非由 Test 搜尋得到。

集成與 FS 將 2014 年以種子 42、依事件標籤分層等分為互不重疊的 DSEL 與 Calibration。DSEL 估計動態方法的局部分類能力；Calibration 選擇靜態成員與分類門檻，未另執行機率校準。兩者均不加入基礎模型擬合，亦無年內先後順序。Baseline 則合用這兩部分，以 2014 年全部 3,504 筆紀錄選擇超參數與門檻；三種訓練策略的擬合資料與集成一致。

\input{generated/event-results-20261103/event-partitions.tex}

2014 年驗證資料共 33 筆事件，分至 DSEL 與 Calibration 後各為 17 與 16 筆，故候選排名與門檻可能受少量事件影響。所有方法固定後才評估 2015--2018 年，各測試年度之間不重新訓練。各年度紀錄數另列於附錄~\ref{sec:appendix-data-fs}。

圖~\ref{fig:method-time-window} 對照一般視窗前移與本次固定配置；上半部僅為概念示意，多輪更新並非本次已執行的實驗。

\begin{figure}[H]
\centering
\includegraphics[width=\linewidth]{figures/time-window-concept.png}
\caption{通用滑動視窗概念與本研究固定年份配置}\label{fig:method-time-window}
\par\smallskip{\footnotesize (資料來源：本研究)\par 上半部為時間配置示意，Val. 表示驗證資料；區段寬度不表示本研究的年份數。下半部之 DSEL 與 Calibration 均來自 2014 年，兩者沒有年內先後順序。\par}
\end{figure}

\subsection{共同前處理、採樣與隨機種子}\label{sec:method-common}

兩項實驗均先以 Old-fit 擬合平均數補值與標準化參數，再將相同轉換套用於 New-fit、驗證與 Test，以維持共同特徵表示。Baseline 的 Retrain 合併轉換後的 Old-fit 與 New-fit。採樣只作用於各模型所屬的擬合資料，驗證與 Test 保留原始類別分布；實驗二另於採樣前執行特徵選取。

表~\ref{tab:method-common-settings} 列出共同設定。Baseline 包含四種採樣條件，集成模型池使用其中 Under、Over 與 Hybrid。第二章已說明各採樣方法的作用，本章沿用相同名稱。

\Needspace{11\baselineskip}
\begingroup
\small\setstretch{1.12}\setlength{\tabcolsep}{4pt}
\begin{longtable}{@{}>{\raggedright\arraybackslash}p{0.18\linewidth}>{\raggedright\arraybackslash}p{0.74\linewidth}@{}}
\caption{共同前處理與採樣設定}\label{tab:method-common-settings}\\
\toprule
項目 & 設定 \\
\midrule\endfirsthead
\multicolumn{2}{c}{表~\thetable（續）}\\\toprule
項目 & 設定 \\
\midrule\endhead
\midrule\multicolumn{2}{r}{續下頁}\\\endfoot
\bottomrule\endlastfoot
補值與縮放 & Old-fit 平均數補值，保留全空欄位的輸入維度；以補值後的 Old-fit 擬合 StandardScaler \\
None & 不採樣，保留擬合資料的類別分布 \\
Under & TomekLinks；\nolinkurl{sampling_strategy=auto} \\
Over & ADASYN；\nolinkurl{sampling_strategy=auto}、\nolinkurl{n_neighbors=5} \\
Hybrid & SMOTEENN；SMOTE 的 \nolinkurl{k_neighbors=5}；ENN 的 \nolinkurl{n_neighbors=3}、\nolinkurl{kind_sel=all} \\
\end{longtable}
\noindent{\footnotesize (資料來源：本研究實驗程式與固定版本套件設定整理)}
\endgroup

模型、採樣及具有隨機性的特徵選取程序，使用 42、123、456、789、2024、7、21、84、168、999 共十個種子，可接收種子的元件依各次設定傳遞；2014 年的資料列分配則固定，不隨模型種子或 FS 條件改變。軟體環境與補充訓練控制見附錄~\ref{app:experimental-details}。十種子規則適用於已完成的事件實驗；深度模型歷史結果另見第~\ref{sec:method-phase1}~節的「深度模型的歷史設定」。

\clearpage

\section{實驗一：Baseline 與集成分類器比較}\label{sec:method-experiment1}

本實驗包含單一分類器基準與模型池集成兩部分，前者比較資料期間、採樣及分類器，後者比較共同模型池的整合規則。

\subsection{Baseline 分類器比較}\label{sec:method-phase1}

Baseline 流程如圖~\ref{fig:method-baseline}，候選包含第二章介紹的 XGBoost、LightGBM、Random Forest、Logistic Regression、TabM 與 TabICL。四種傳統分類器已完成事件標籤實驗；兩種深度模型目前先列歷史結果，與事件結果分表呈現。

\begin{figure}[htbp]
\centering
\includegraphics[width=\linewidth]{figures/experiment1-baseline.png}
\caption{實驗一之 Baseline 分類器比較架構}\label{fig:method-baseline}
\par\smallskip{\footnotesize (資料來源：本研究)\par}
\end{figure}

\subsubsection{訓練策略與執行流程}

各分類器分別建立三種模型：Old 使用 Old-fit，New 使用 New-fit，Retrain 合併兩者從頭訓練。New 亦為重新訓練，未延續 Old 模型權重。各策略搭配表~\ref{tab:method-common-settings} 的四種採樣，採用第~\ref{sec:method-common}~節的共同前處理。

對每個「分類器 \ensuremath{\times} 訓練策略 \ensuremath{\times} 採樣」組合，先依驗證 ROC-AUC 選出超參數，再以十個種子重新訓練，並分別依驗證 F1-score 決定門檻，最後計算 Test 指標。已完成部分共 \(4\times3\times4=48\) 個組合；深度模型的歷史設定不套用此完成範圍。

\subsubsection{超參數候選與選擇}\label{sec:method-phase1-params}

每個組合評估 24 組候選：候選 0 為目前集成測試採用的預設參數，其餘 23 組依固定亂數從表~\ref{tab:method-ml-search} 抽出，不窮舉全部排列。候選均以種子 42 擬合，驗證 AUC 相同時取候選編號較小者。選定參數後的十種子評估沿用前述流程，不依 Test 改選。

\begingroup
\footnotesize
\setstretch{1.1}
\setlength{\tabcolsep}{4pt}
\renewcommand{\arraystretch}{1.15}
\setlength{\LTleft}{0pt plus 1fill}
\setlength{\LTright}{0pt plus 1fill}
\begin{longtable}{@{}>{\raggedright\arraybackslash}p{\dimexpr 0.2100\linewidth-0.8400\tabcolsep\relax}>{\raggedright\arraybackslash}p{\dimexpr 0.3100\linewidth-1.2400\tabcolsep\relax}>{\raggedright\arraybackslash}p{\dimexpr 0.4800\linewidth-1.9200\tabcolsep\relax}@{}}
\caption{Baseline 傳統分類器之超參數候選集合}\label{tab:method-ml-search}\\
\toprule
\textbf{分類器} & \textbf{參數項目} & \textbf{超參數候選值} \\
\midrule
\endfirsthead
\multicolumn{3}{c}{表~\thetable（續）}\\
\toprule
\textbf{分類器} & \textbf{參數項目} & \textbf{超參數候選值} \\
\midrule
\endhead
\midrule
\multicolumn{3}{r}{續下頁}\\
\endfoot
\bottomrule
\endlastfoot
\multirow[t]{9}{=}{XGBoost} & n\_\allowbreak{}estimators & 100、300、500、800、1000 \\*
\addlinespace[1pt]
 & learning\_\allowbreak{}rate & 0.01、0.03、0.05、0.1、0.2 \\*
\addlinespace[1pt]
 & max\_\allowbreak{}depth & 3、5、7、10 \\*
\addlinespace[1pt]
 & min\_\allowbreak{}child\_\allowbreak{}weight & 1、3、5、20、50 \\*
\addlinespace[1pt]
 & subsample & 0.7、0.85、1.0 \\*
\addlinespace[1pt]
 & colsample\_\allowbreak{}bytree & 0.7、0.85、1.0 \\*
\addlinespace[1pt]
 & gamma & 0.0、0.1、0.2 \\*
\addlinespace[1pt]
 & reg\_\allowbreak{}alpha & 0.0、0.1、1.0 \\*
\addlinespace[1pt]
 & reg\_\allowbreak{}lambda & 0.5、1.0、3.0、10.0 \\
\addlinespace[1pt]
\multirow[t]{9}{=}{LightGBM} & num\_\allowbreak{}leaves & 31、50、100、150 \\*
\addlinespace[1pt]
 & max\_\allowbreak{}depth & 5、7、10、15、-1 \\*
\addlinespace[1pt]
 & learning\_\allowbreak{}rate & 0.01、0.03、0.05、0.1、0.2 \\*
\addlinespace[1pt]
 & n\_\allowbreak{}estimators & 100、200、500、800、1000 \\*
\addlinespace[1pt]
 & min\_\allowbreak{}child\_\allowbreak{}samples & 20、50、100 \\*
\addlinespace[1pt]
 & subsample & 0.7、0.85、1.0 \\*
\addlinespace[1pt]
 & colsample\_\allowbreak{}bytree & 0.7、0.85、1.0 \\*
\addlinespace[1pt]
 & reg\_\allowbreak{}alpha & 0.0、0.1、1.0 \\*
\addlinespace[1pt]
 & reg\_\allowbreak{}lambda & 0.0、0.1、1.0、3.0 \\
\addlinespace[1pt]
\multirow[t]{5}{=}{Random Forest} & n\_\allowbreak{}estimators & 200、400、600、800 \\*
\addlinespace[1pt]
 & max\_\allowbreak{}depth & None、8、12、16、20、24、32 \\*
\addlinespace[1pt]
 & min\_\allowbreak{}samples\_\allowbreak{}leaf & 1、2、4、8、16 \\*
\addlinespace[1pt]
 & min\_\allowbreak{}samples\_\allowbreak{}split & 2、4、8、16、32 \\*
\addlinespace[1pt]
 & max\_\allowbreak{}features & sqrt、log2、0.2、0.35、0.5、0.65 \\
\addlinespace[1pt]
\multirow[t]{3}{=}{Logistic Regression} & C & 0.01、0.1、1.0、10.0、100.0 \\*
\addlinespace[1pt]
 & solver & lbfgs、saga \\*
\addlinespace[1pt]
 & max\_\allowbreak{}iter & 500、1000、2000 \\
\end{longtable}
\noindent{\footnotesize (資料來源：本研究整理)}
\endgroup

表中 None 與深度值 -1 均表示不設樹深度上限；Random Forest 的小數型 \nolinkurl{max_features} 表示每次分裂考量的特徵比例。各組合選中的候選編號保存於執行紀錄。

\subsubsection{深度模型的歷史設定}\label{sec:method-deep-reference}

TabM 的既有紀錄保存了驗證選定的網路寬度、區塊數、成員數及最佳化參數。表~\ref{tab:deep-reference-settings} 展示 Old 1999--2011 年、Hybrid 採樣的一組設定，並非所有情境共用的參數，展示情境亦未依測試分數挑選。TabICL 的既有入口支援不同套件後端，歷史 CSV 不足以確認實際後端及完整選參過程，故其模型版本仍待核對。

\input{generated/deep-baseline-reference-20261103/settings.tex}

兩者歷史結果採公司最終狀態標籤：Old／New 分界在 1999--2014 年間移動，共 15 組重疊設定，Test 固定為 2015--2018 年，驗證方式與本次事件協定不同。第四章暫列既有指標；完成事件標籤及共同驗證規則的重跑後，才能更新為六分類器的同條件比較。

\subsubsection{集成基礎模型的選擇與完成範圍}\label{sec:method-selection-status}

研究框架以共同驗證 ROC-AUC 選定基礎分類器及參數，再建立 Old／New 六模型池。第四章先彙整四種傳統分類器在六個池位置（Old／New \ensuremath{\times} Under、Over、Hybrid）的驗證表現；正式選模仍待完整候選齊備。現有集成與 FS 結果來自 XGBoost 預設參數（候選 0）的流程測試，尚未依完整六分類器驗證結果重建模型池。以下交代該測試的可重現設定，並以此界定第四章現有結果的範圍。

\subsection{Old／New 集成分類器比較}\label{sec:method-phase2}

集成流程如圖~\ref{fig:method-ensemble}。各方法共用同一基礎分類器及六模型池，差異限於成員選擇與整合規則；模型選擇與完成範圍依前述說明。

\begin{figure}[htbp]
\centering
\includegraphics[width=\linewidth]{figures/experiment1-ensemble.png}
\caption{實驗一之 Old／New 集成分類器架構}\label{fig:method-ensemble}
\par\smallskip{\footnotesize (資料來源：本研究)\par}
\par\smallskip{\footnotesize 圖中 Validation 為驗證用途的概括表示；實際分為互不重疊的 DSEL 與 Calibration。前者估計局部分類能力，後者選擇靜態候選及分類門檻。\par}
\end{figure}

\subsubsection{模型池建立與分類器設定}

依共同前處理完成轉換後，Old-fit 與 New-fit 各自採用三種採樣並訓練分類器，形成表~\ref{tab:method-pool} 的六個位置。實驗一使用完整 18 個特徵，基礎分類器不另加入類別權重。

\Needspace{11\baselineskip}
\begingroup
\small
\setstretch{1.15}
\setlength{\tabcolsep}{4pt}
\renewcommand{\arraystretch}{1.25}
\setlength{\LTleft}{0pt plus 1fill}
\setlength{\LTright}{0pt plus 1fill}
\begin{longtable}{@{}>{\raggedright\arraybackslash}p{\dimexpr 0.28\linewidth-1.1200\tabcolsep\relax}>{\raggedright\arraybackslash}p{\dimexpr 0.25\linewidth-1.0000\tabcolsep\relax}>{\raggedright\arraybackslash}p{\dimexpr 0.47\linewidth-1.8800\tabcolsep\relax}@{}}
\caption{Old／New 六模型池組成}\label{tab:method-pool}\\
\toprule
\textbf{模型名稱} & \textbf{擬合資料} & \textbf{採樣方法} \\
\midrule
\endfirsthead
\multicolumn{3}{c}{表~\thetable（續）}\\
\toprule
\textbf{模型名稱} & \textbf{擬合資料} & \textbf{採樣方法} \\
\midrule
\endhead
\midrule
\multicolumn{3}{r}{續下頁}\\
\endfoot
\bottomrule
\endlastfoot
Old\_under & Old-fit & TomekLinks \\
\addlinespace[2pt]
Old\_over & Old-fit & ADASYN \\
\addlinespace[2pt]
Old\_hybrid & Old-fit & SMOTEENN \\
\addlinespace[2pt]
New\_under & New-fit & TomekLinks \\
\addlinespace[2pt]
New\_over & New-fit & ADASYN \\
\addlinespace[2pt]
New\_hybrid & New-fit & SMOTEENN \\*
\end{longtable}
\noindent{\footnotesize (資料來源：本研究整理)}
\endgroup

目前測試的 XGBoost 設定為 \nolinkurl{objective=binary:logistic}、\nolinkurl{eval_metric=auc}、\nolinkurl{tree_method=hist}，其餘依固定版本預設值；未另搜尋池內超參數，亦未啟用早停。Baseline 與集成的參數及驗證用途不同，兩者分數差不單獨解釋為集成效果。Retrain 僅列於 Baseline，不納入此模型池。

\subsubsection{靜態成員與機率整合}

沿用第~\ref{sec:method-framework}~節的 \(\mathcal{O}\) 與 \(\mathcal{N}\)，跨期間混合候選須滿足：

\begin{equation}\label{eq:method-subsets}
S\subseteq\mathcal{O}\cup\mathcal{N},\qquad
S\cap\mathcal{O}\ne\varnothing,\qquad
S\cap\mathcal{N}\ne\varnothing.
\end{equation}

兩池各有三個模型，故混合候選共 \((2^3-1)(2^3-1)=49\) 種，依成員數分組如表~\ref{tab:method-combinations}。

\begingroup
\small
\setstretch{1.15}
\setlength{\tabcolsep}{4pt}
\renewcommand{\arraystretch}{1.25}
\setlength{\LTleft}{0pt plus 1fill}
\setlength{\LTright}{0pt plus 1fill}
\begin{longtable}{@{}>{\raggedright\arraybackslash}p{\dimexpr 0.55\linewidth-1.1000\tabcolsep\relax}>{\raggedright\arraybackslash}p{\dimexpr 0.45\linewidth-0.9000\tabcolsep\relax}@{}}
\caption{不同模型數量的 Old／New 混合靜態組合}\label{tab:method-combinations}\\
\toprule
\textbf{集成模型數量} & \textbf{混合組合數} \\
\midrule
\endfirsthead
\multicolumn{2}{c}{表~\thetable（續）}\\
\toprule
\textbf{集成模型數量} & \textbf{混合組合數} \\
\midrule
\endhead
\midrule
\multicolumn{2}{r}{續下頁}\\
\endfoot
\bottomrule
\endlastfoot
2 & 9 \\
\addlinespace[2pt]
3 & 18 \\
\addlinespace[2pt]
4 & 15 \\
\addlinespace[2pt]
5 & 6 \\
\addlinespace[2pt]
6 & 1 \\
\addlinespace[2pt]
合計 & 49 \\*
\end{longtable}
\noindent{\footnotesize (資料來源：本研究整理)}
\endgroup

對任一候選子集 \(S\)，以等權平均形成集成機率：

\begin{equation}\label{eq:method-probability}
\hat p_S(x)=\frac{1}{|S|}\sum_{i\in S}\hat p_i(x),
\end{equation}

其中，\(\hat p_i(x)\) 為第 \(i\) 個分類器對樣本 \(x\) 的正類機率，\(|S|\) 為成員數。另保留六個單模型，以及單一期間控制組 Old3、New3。完整六模型等權集成稱為 Equal6，已包含於 49 種混合候選；Old3 與 New3 則不在混合候選範圍內。

\subsubsection{動態選擇的實作設定}

依第二章所述機制，以完整六模型池比較 DCS 的 OLA、LCA，以及 DES 的 KNORA-E、KNORA-U。四種方法共用 DSEL，使用 DESlib 0.3.7，鄰居數固定為 \(k=7\)。DCS 採 \nolinkurl{best} 規則，局部能力平手時依固定模型順序選擇；DES 採 \nolinkurl{soft} 機率整合。額外動態剪枝與樣本難度分流均關閉，即 \nolinkurl{DFP=False}、\nolinkurl{with_IH=False}。

所有基礎分類器對某樣本判斷一致時，套件直接平均機率；因此 DCS 的輸出不一定每次都來自單一模型。其餘局部選擇與權重處理沿用該版本標準實作\autocite{deslib-ola-docs,deslib-lca-docs,deslib-knorae-docs,deslib-knorau-docs}，未另修改演算法。

\subsubsection{靜態候選與分類門檻的決定}

先在 Calibration 對 49 種混合候選計算 AP，選出最高者，記為 V-static（執行紀錄名稱為 ValidationSelectedStatic）；平手時依固定方法名稱的字典序決定。全部候選的 Test 表現僅供組成分析，不用來回選成員。

各單模型、控制組、靜態候選及動態方法，均在 Calibration 搜尋 F1-score 最高的門檻：

\begin{equation}\label{eq:method-threshold}
\mathcal{T}=\{0.01,0.02,\ldots,0.99\},\qquad
\tau_m^*=\underset{\tau\in\mathcal{T}}{\arg\max}\;
F1_{\mathrm{Calibration}}(m,\tau).
\end{equation}

平手取最小門檻，以 \(\hat p_m(x)\ge\tau_m^*\) 判為正類；V-static 沿用所選候選的門檻。成員與門檻固定後，才計算 Test 的整體與逐年指標。

\clearpage

\section{實驗二：特徵選取集成分類器比較}\label{sec:method-phase3}

實驗二沿用實驗一的集成設定，只在採樣前加入特徵選取，再重新訓練模型池，以檢查特徵精簡對分類表現的影響。流程如圖~\ref{fig:method-fs}。

\begin{figure}[htbp]
\centering
\includegraphics[width=\linewidth]{figures/experiment2-fs-ensemble.png}
\caption{實驗二之特徵選取集成分類器架構}\label{fig:method-fs}
\par\smallskip{\footnotesize (資料來源：本研究)\par}
\par\smallskip{\footnotesize 各 FS 方法平行執行，並在選取後分別重建模型池與動態選擇器。\par}
\end{figure}

\subsection{執行流程與特徵條件}

處理順序為「資料分區 \ensuremath{\rightarrow} 前處理 \ensuremath{\rightarrow} 特徵選取 \ensuremath{\rightarrow} 採樣 \ensuremath{\rightarrow} 模型訓練」。各選取器僅擬合轉換後、尚未採樣的 Old-fit；所選欄位及順序固定後，一致套用於 New-fit、DSEL、Calibration 與 Test。每個條件重新建立六模型池及動態選擇器，使分類器與局部鄰近區域使用相同輸入空間。

MI、CART、RFE 與 SHAP 排序分別獨立執行，不串接使用。各方法設定名義保留比例 \(r=0.5\) 與 \(r=0.8\)，依下式決定保留數：

\begin{equation}\label{eq:method-fs-count}
k_{\mathrm{FS}}=\max\{1,\lfloor d\times r\rfloor\}.
\end{equation}

原始特徵數 \(d=18\)，兩設定分別保留 9 與 14 個；加上完整特徵的 no-FS，共九種條件。

兩個比例用於比較較大與較小幅度的刪減，並非由驗證或 Test 搜尋的最佳比例；80\% 為名義值，實際數量依向下取整計算。

\subsection{選取依據與固定參數}\label{sec:method-fs-parameters}

表~\ref{tab:method-fs} 摘要選取依據，表~\ref{tab:method-fs-parameters} 列出固定參數；方法原理與 XAI 定位沿用第二章。

\begingroup
\small
\setstretch{1.15}
\setlength{\tabcolsep}{4pt}
\renewcommand{\arraystretch}{1.25}
\setlength{\LTleft}{0pt plus 1fill}
\setlength{\LTright}{0pt plus 1fill}
\begin{longtable}{@{}>{\raggedright\arraybackslash}p{\dimexpr 0.15\linewidth-0.6000\tabcolsep\relax}>{\raggedright\arraybackslash}p{\dimexpr 0.28\linewidth-1.1200\tabcolsep\relax}>{\raggedright\arraybackslash}p{\dimexpr 0.57\linewidth-2.2800\tabcolsep\relax}@{}}
\caption{特徵選取方法與操作方式}\label{tab:method-fs}\\
\toprule
\textbf{方法} & \textbf{方法定位} & \textbf{特徵選取依據} \\
\midrule
\endfirsthead
\multicolumn{3}{c}{表~\thetable（續）}\\
\toprule
\textbf{方法} & \textbf{方法定位} & \textbf{特徵選取依據} \\
\midrule
\endhead
\midrule
\multicolumn{3}{r}{續下頁}\\
\endfoot
\bottomrule
\endlastfoot
MI & Filter & 計算各特徵與分類標籤的互資訊，依分數選取前 \(k_{\mathrm{FS}}\) 個特徵 \\
\addlinespace[2pt]
CART & 模型重要性篩選，對應 Embedded 類 & 以決策樹學得的特徵重要性排序，選取前 \(k_{\mathrm{FS}}\) 個特徵 \\
\addlinespace[2pt]
RFE & Wrapper & 以決策樹反覆擬合並移除特徵，直至保留指定數量 \\
\addlinespace[2pt]
SHAP 排序 & 模型解釋後的特徵篩選 & 以 LightGBM 計算訓練樣本的樹模型貢獻值，依平均絕對值排序選取特徵 \\*
\end{longtable}
\noindent{\footnotesize (資料來源：本研究整理)}
\endgroup

MI 與 RFE 由 scikit-learn 實作\autocite{sklearn-fs-docs}；SHAP 排序使用 LightGBM 的 \nolinkurl{pred_contrib=True} 輸出，排除基準值欄位後計算平均絕對貢獻\autocite{lightgbm-booster-docs}。CART、RFE 與 SHAP 的平衡類別權重只作用於選取器，不改變後續基礎分類器設定。未另指定的參數沿用表列套件版本預設值，亦不使用 Calibration 或 Test 調整選取器。

\begingroup
\footnotesize\setstretch{1.12}\setlength{\tabcolsep}{3pt}
\begin{longtable}{@{}>{\raggedright\arraybackslash}p{0.13\linewidth}>{\raggedright\arraybackslash}p{0.31\linewidth}>{\raggedright\arraybackslash}p{0.49\linewidth}@{}}
\caption{實驗二特徵選取之主要參數與執行規則}\label{tab:method-fs-parameters}\\
\toprule
方法 & 參數或項目 & 實際設定 \\
\midrule\endfirsthead
\multicolumn{3}{c}{表~\thetable（續）}\\\toprule
方法 & 參數或項目 & 實際設定 \\
\midrule\endhead
\midrule\multicolumn{3}{r}{續下頁}\\\endfoot
\bottomrule\endlastfoot
\multirow[t]{2}{=}{共同設定} & 擬合資料 & 前處理後的 Old-fit；於採樣前執行 \\*
 & 保留數與隨機種子 & 9、14 個；\nolinkurl{random_state} 隨十個模型種子設定 \\
\multirow[t]{3}{=}{MI} & 函式與連續特徵 & \nolinkurl{mutual_info_classif}；\nolinkurl{discrete_features=auto}，本研究密集數值矩陣按連續特徵處理 \\*
 & 鄰居數 & \nolinkurl{n_neighbors=3}（套件預設） \\*
 & 選取規則 & 互資訊分數由高至低；平手保留原欄位先後順序 \\
\multirow[t]{5}{=}{CART} & 估計器 & \nolinkurl{DecisionTreeClassifier}；\nolinkurl{class_weight=balanced} \\*
 & 分裂規則 & \nolinkurl{criterion=gini}、\nolinkurl{splitter=best}（預設） \\*
 & 樹的大小 & \nolinkurl{max_depth=None}、\nolinkurl{min_samples_split=2}、\nolinkurl{min_samples_leaf=1}（預設） \\*
 & 剪枝與特徵抽樣 & \nolinkurl{ccp_alpha=0}、\nolinkurl{max_features=None}（預設） \\*
 & 選取規則 & 依特徵重要性排序；平手保留原欄位先後順序 \\
\multirow[t]{3}{=}{RFE} & 估計器 & 與 CART 相同的平衡類別權重決策樹 \\*
 & 移除步長及停止條件 & \nolinkurl{step=1}；\nolinkurl{n_features_to_select} 分別為 9、14 \\*
 & 輸出順序 & 依選取遮罩保留原欄位順序 \\
\multirow[t]{6}{=}{SHAP 排序} & 估計器 & \nolinkurl{LGBMClassifier}；\nolinkurl{n_estimators=100}、\nolinkurl{class_weight=balanced} \\*
 & 樹與學習率 & \nolinkurl{boosting_type=gbdt}、\nolinkurl{num_leaves=31}、\nolinkurl{max_depth=-1}、\nolinkurl{learning_rate=0.1}（預設） \\*
 & 葉節點與抽樣 & \nolinkurl{min_child_samples=20}、\nolinkurl{colsample_bytree=1.0}、\nolinkurl{subsample=1.0}（預設） \\*
 & 正則化與執行緒 & \nolinkurl{reg_alpha=0}、\nolinkurl{reg_lambda=0}（預設）；\nolinkurl{n_jobs=4} \\*
 & 重要性計算 & 在 Old-fit 使用 \nolinkurl{pred_contrib=True}；排除基準值欄位，再計算平均絕對貢獻值 \\*
 & 選取規則 & 依貢獻值排序；平手保留原欄位先後順序 \\
\end{longtable}
\noindent{\footnotesize (資料來源：本研究特徵選取程式與固定版本套件設定整理)\par scikit-learn 1.8.0；LightGBM 4.6.0。未列項目沿用相同版本預設值。}
\endgroup

\subsection{配對比較方式}

各條件沿用 49 種混合靜態候選、Old3、New3、Equal6 與四種動態方法，使用相同資料列分區、模型種子及 Test 樣本。Calibration 在各條件下重新選擇靜態候選與分類門檻。因此 V-static 的差值包含特徵與驗證選定成員的改變；其他方法維持模型池位置及整合規則，仍依各條件重新訓練。動態方法對個別樣本選出的成員也可能隨特徵空間改變。

對固定方法 \(m\)、指標 \(M\)、種子 \(s\) 與特徵條件 \(f\)，配對差值為：

\begin{equation}\label{eq:method-delta}
\Delta M_{s,m,f}=M_{s,m,f}-M_{s,m,\mathrm{no\_fs}}.
\end{equation}

正值表示高於 no-FS。即使成員組合固定，F1-score 與 Recall 差異仍包含重新訓練與各自門檻選擇的影響。各次所選欄位名稱、順序與數量均保存，十種子的入選次數列於附錄~\ref{sec:appendix-data-fs}，僅作一致性的描述，未另作排名顯著性檢定。

\clearpage
\section{實驗驗證準則}\label{sec:method-validation}

本節統一定義指標、結果彙整與統計比較。資料用途與選擇程序依前述設定執行，驗證最高分不作為最終預測表現。

\subsection{分類效能評估指標}

令關注類別為正類。TP 與 TN 分別為正確分類的正類與負類紀錄數；FN 為正類被判為負類的紀錄數；FP 為負類被判為正類的紀錄數。實證案例的正類對應第~\ref{sec:method-data}~節定義的標籤。

\textbf{一、AUC}

ROC 曲線以真陽性率（True Positive Rate, TPR）為縱軸，假陽性率（False Positive Rate, FPR）為橫軸，定義為：

\begin{equation}\label{eq:method-roc}
\mathrm{TPR}=\frac{TP}{TP+FN},\qquad
\mathrm{FPR}=\frac{FP}{FP+TN}.
\end{equation}

AUC 為 ROC 曲線下的面積，由預測機率計算，不使用分類門檻。

\textbf{二、F1-score 與 Recall}

Precision 表示預測為正類的紀錄中，實際正類者所占比例；Recall 表示實際正類紀錄中被正確辨識的比例。兩者與 F1-score 定義如下：

\begin{equation}\label{eq:recall}
\mathrm{Precision}=\frac{TP}{TP+FP},\qquad
\mathrm{Recall}=\frac{TP}{TP+FN},
\end{equation}

\begin{equation}\label{eq:method-f1}
F1=\frac{2\,\mathrm{Precision}\,\mathrm{Recall}}
{\mathrm{Precision}+\mathrm{Recall}}
=\frac{2TP}{2TP+FP+FN}.
\end{equation}

Recall 亦即 TPR。兩項指標使用各方法已固定的驗證門檻；Baseline 使用完整 2014 年，集成與 FS 使用 Calibration。遇到未定義的零分母時，統一記為零。

\Needspace{6\baselineskip}
\textbf{三、Average Precision}

AP 用於實驗一之集成比較與實驗二的靜態候選選擇，依不同門檻的 Precision 與 Recall 計算：

\begin{equation}\label{eq:method-ap}
\mathrm{AP}=\sum_n(R_n-R_{n-1})P_n,
\end{equation}

其中，\(P_n\) 與 \(R_n\) 為第 \(n\) 個操作點的 Precision 與 Recall。本研究依 scikit-learn 定義計算 AP\autocite{sklearn-ap-docs}，用於驗證選模；Test 主文指標為 AUC、F1-score 與 Recall。

\textbf{四、特徵精簡率}

除預測指標外，實驗二也記錄刪除了多少特徵。令原始特徵數為 \(d\)、保留數為 \(k_{\mathrm{FS}}\)，特徵精簡率定義為：
\begin{equation}\label{eq:method-reduction}
R_{\mathrm{FS}}=\left(1-\frac{k_{\mathrm{FS}}}{d}\right)\times100\%.
\end{equation}
保留 9 與 14 個特徵時，精簡率分別為 50.00\% 與 22.22\%；no-FS 為 0\%。此指標只描述輸入維度的減少。

\subsection{結果彙整}

整體結果先合併 2015--2018 年的測試標籤與預測，再計算指標，稱為合併四年結果（pooled），不等同各年度分數的平均。逐年結果另外計算，模型與門檻均沿用驗證決策。相同設定的十種子結果以平均值及標準差呈現；種子共用同一份 Test，標準差僅描述執行變異，不代表獨立資料的不確定性。歷史深度模型依 Baseline 節所述的重疊設定彙整，不併入十種子比較。

\subsection{以公司為單位的統計比較}\label{sec:method-statistics}

本研究補作公司群聚配對 bootstrap，以同一批重新抽取的測試公司比較方法差值。事件標籤為主要分析，最終狀態標籤另作對照，各自檢定與校正。此為既有實驗完成後制定的事後探索性分析，沿用保存的模型、特徵、成員與門檻，不因檢定結果重新選模。統計原理及群聚相關性的文獻依據已於第二章說明，以下列出實際估計量與重抽程序。

\subsubsection{比較對象與估計量}

實驗一以完整六模型等權集成 Equal6 為共同基準，分別比較 Old3、New3、V-static、DCS-OLA、DCS-LCA、KNORA-E 及 KNORA-U，共七組比較。Baseline 因參數及驗證用途不同，不納入此檢定。

實驗二將 MI、CART、SHAP 與 RFE 的兩種保留比例，各自與相同方法的 no-FS 結果配對。比較方法為 Equal6、V-static、DCS-OLA、DCS-LCA、KNORA-E 與 KNORA-U，因此共有 \(8\times6=48\) 組比較。V-static 與固定方法的比較界線依實驗二所述規則處理。

每組比較均計算 AUC、F1-score 及 Recall，合計 \((7+48)\times3=165\) 項檢定。若 \(M\) 為任一指標，\(S=10\) 為種子數，方法 A 相對方法 B 的差值定義為：
\begin{equation}\label{eq:stats-difference}
\widehat{\Delta}_M=\frac{1}{S}\sum_{s=1}^{S}
\left[M(y,p_{A,s};\tau_{A,s})-M(y,p_{B,s};\tau_{B,s})\right].
\end{equation}
其中，\(p\) 為模型預測機率，\(\tau\) 為驗證程序固定的門檻；AUC 不使用門檻。此估計量先計算各種子的 pooled 指標，再平均其差值，不將十個種子的機率平均後另建一個預測模型。

\subsubsection{重抽、信賴區間與多重比較}

每次 bootstrap 從 Test 所含的 \(C\) 家公司中等機率、有放回抽取 \(C\) 次。公司被抽中後，保留該公司在 2015--2018 年的全部測試紀錄，並將同一組公司抽取次數同時套用於所有種子、方法及 FS 設定。如此可保留同公司跨年的相依性及方法間的配對關係；各次重抽均重新計算指標，但不重新訓練或調整門檻。

事件標籤分析使用 9,999 次重抽，最終狀態標籤對照使用 49,999 次，兩批皆保留既有分析規格，未因顯著性結果改變重抽次數。為使統計分析可重現，bootstrap 重抽樣的隨機種子皆設為 20261027。由差值重抽分布的第 2.5 與第 97.5 百分位數建立逐項 95\% 信賴區間。以 \(\Delta_M^{*(b)}\) 表示第 \(b\) 次重抽差值，針對虛無假設 \(H_0:\Delta_M=0\)，以置中重抽誤差分布計算雙尾近似 \(p\) 值：
\begin{equation}\label{eq:stats-bootstrap-p}
p_M=\frac{1+\sum_{b=1}^{B}\mathbf{1}\!\left\{
\left|\Delta_M^{*(b)}-\widehat{\Delta}_M\right|\geq
\left|\widehat{\Delta}_M\right|\right\}}{B+1},
\qquad B\in\{9{,}999,49{,}999\}.
\end{equation}
式中加一避免報告零機率；兩種標籤的最小可報值依序為 0.0001 與 0.00002。有限次重抽仍有 Monte Carlo 誤差，尾端重抽次數很少時，校正後判斷可能對此敏感。

每種標籤下的兩項實驗共 165 個 \(p\) 值一併使用 Holm 法\autocite{holm1979multiple} 校正。校正後數值記為 \(p_H\)，小於 0.05 時判為達到統計顯著。表中的 95\% 信賴區間各自計算，沒有作同樣的多重比較校正，所以顯著性仍以 \(p_H\) 為準。未達顯著只表示本次證據不足，不能據此認定兩方法相同。

此分析未涵蓋訓練資料變動及跨公司產業、景氣關聯。事件檢定摘要見第~\ref{sec:results-event-statistics}~節，完整比較見附錄。
\endgroup
